\documentclass[12pt]{article}
\usepackage[utf8]{inputenc}
\usepackage[T1]{fontenc}
\usepackage{hyperref}
\title{Yang Zhuang Language Resources for K-12 Curriculum Development}
\author{Compiled by Claude Sonnet 5.5 (claude-sonnet-5-5) for Andr\'as Kornai}
\date{}
\begin{document}
\maketitle

\section{Introduction}
Yang Zhuang (ISO 639-3 zyg, Glottocode yang1286) is a Tai language of southwestern Guangxi, China. Wikipedia places it in Napo, Jingxi and Debao counties and gives about 770,000 speakers (2004). Joshua Project gives a higher figure. Alternate names reported include Dejing Zhuang, Jingxi Zhuang and Nung Giang, and a related community is reported in Vietnam. The search found no school curricula or textbooks specific to this language, so little K-12 material is documented here.

\section{Primary Educational Resources}
No primers, textbooks or school programmes specific to Yang Zhuang turned up in the search results. This section is therefore empty.

\section{Linguistic and Reference Materials}
\begin{itemize}
\item Wikipedia article on Yang Zhuang: \url{https://en.wikipedia.org/wiki/Yang_Zhuang_language}
\item Omniglot page on Yang Zhuang, covering dialects and the traditional Chinese-character-based writing: \url{https://omniglot.com/writing/yangzhuang.htm}
\item Glottolog language entry (ISO zyg): \url{https://glottolog.org/resource/languoid/iso/zyg}
\item Wikitongues language page: \url{https://wikitongues.org/languages/zyg/}
\end{itemize}

\section{Religious and Community Resources}
\begin{itemize}
\item Joshua Project profile of Yang Zhuang in China: \url{https://www.joshuaproject.net/index.php/people_groups/19686/CH}
\item Joshua Project profile of Yang Zhuang in Vietnam: \url{https://www.joshuaproject.net/index.php/people_groups/19686/VM}
\end{itemize}

\section{Digital and Audio Resources}
\begin{itemize}
\item Global Recordings Network page, ``Zhuang, Yang'': \url{https://globalrecordings.net/en/language/zyg}
\end{itemize}

\section{Curriculum Development Context}
Omniglot reports no standard orthography for Yang Zhuang, and says stories have traditionally been written in a Chinese-character-based script. Glottolog classifies the language as vigorous. Any curriculum would first need an orthography decision, and this report found no existing teaching materials.

\end{document}